\documentclass[border=10pt]{standalone}
\usepackage[no-math]{fontspec}
\usepackage[brazil]{babel}
\usepackage{amsmath,amssymb}
\IfFileExists{C:/Windows/Fonts/Lato_Regular.ttf}{%
  \setmainfont{Lato}[
    Path = C:/Windows/Fonts/,
    UprightFont = Lato_Regular.ttf,
    BoldFont = Lato_Bold.ttf,
    ItalicFont = Lato_Italic.ttf,
    BoldItalicFont = Lato_BoldItalic.ttf
  ]
  \setsansfont{Lato}[
    Path = C:/Windows/Fonts/,
    UprightFont = Lato_Regular.ttf,
    BoldFont = Lato_Bold.ttf,
    ItalicFont = Lato_Italic.ttf,
    BoldItalicFont = Lato_BoldItalic.ttf
  ]
}{%
  \RequirePackage[default]{lato}
}
\renewcommand{\familydefault}{\sfdefault}

\usepackage{xcolor}
\usepackage{tikz}
\usepackage{pgfplots}
\pgfplotsset{compat=1.18}
\usetikzlibrary{arrows.meta,positioning,shapes.geometric,fit,backgrounds,calc,decorations.pathreplacing,patterns}

% Paleta institucional UEFS / Bioengenharia de Solos
\definecolor{bsNavy}{HTML}{1F4268}
\definecolor{bsNavyDark}{HTML}{152E4A}
\definecolor{bsGold}{HTML}{B38A3D}
\definecolor{bsGoldDark}{HTML}{8C6927}
\definecolor{bsTeal}{HTML}{0E7C7B}
\definecolor{bsForest}{HTML}{2D6A4F}
\definecolor{bsAlert}{HTML}{9E2A2B}
\definecolor{bsEarth}{HTML}{8C5A3C}
\definecolor{bsWater}{HTML}{2E86AB}
\definecolor{bsInk}{HTML}{26313D}
\definecolor{bsMuted}{HTML}{5B6573}

\definecolor{bgNavy}{HTML}{EDF2F7}
\definecolor{bgGold}{HTML}{F8F3E9}
\definecolor{bgTeal}{HTML}{EBF6F5}
\definecolor{bgForest}{HTML}{EEF5F0}
\definecolor{bgAlert}{HTML}{FDF2F2}

\tikzset{
  >=Stealth,
  bsBox/.style={
    draw=#1, fill=white, rounded corners=5pt, line width=1.15pt,
    inner sep=7pt, align=center, font=\footnotesize\sffamily, text=bsInk
  },
  bsTint/.style={
    draw=#1, fill=#1!8!white, rounded corners=5pt, line width=1.15pt,
    inner sep=7pt, align=center, font=\footnotesize\sffamily, text=bsInk
  },
  bsBadge/.style={
    fill=#1, text=white, rounded corners=9pt, inner xsep=8pt, inner ysep=3pt,
    font=\scriptsize\bfseries\sffamily
  },
  bsTitle/.style={font=\normalsize\bfseries\sffamily, text=bsNavy, align=center},
  bsSub/.style={font=\footnotesize\sffamily, text=bsMuted, align=center},
  bsArrow/.style={->, >=Stealth, line width=1.35pt},
  bsDArrow/.style={<->, >=Stealth, line width=1.35pt}
}


\begin{document}
\begin{tikzpicture}

  % Título
  \node[bsTitle] at (4.5, 7.25)
    {Dinâmica Fluvial em Curvas Meândricas e Zonas de Atuação da Bioengenharia};
  \node[bsSub] at (4.5, 6.75)
    {Assimetria hidráulica em planta/seção transversal e escalonamento vertical das técnicas na margem};

  % ============================================================
  % PAINEL ESQUERDO: PLANTA DO MEANDRO + SEÇÃO TRANSVERSAL
  % ============================================================
  \node[bsBox=bsNavy!55!white, minimum width=9.6cm, minimum height=8.85cm, anchor=north west] (leftpanel) at (-2.7, 6.10) {};
  \node[bsBadge=bsNavy] at (leftpanel.north) {HIDRODINÂMICA EM CURVA MEÂNDRICA};

  % Planta do meandro
  \begin{scope}[shift={(0.1, 0.45)}]
    \node[font=\scriptsize\bfseries, text=bsMuted, anchor=west] at (-2.45, 5.05) {PLANTA};

    % Caixa de destaque margem côncava (bem abaixo do badge e separada de PLANTA)
    \node[font=\scriptsize\bfseries, text=bsAlert, align=center, fill=bgAlert, draw=bsAlert, rounded corners=3pt, inner xsep=5pt, inner ysep=3pt]
      at (2.05, 4.45) {Margem Côncava (Externa): Alta Velocidade \& Solapamento $\rightarrow$ Proteção de Base};

    % Canal meândrico
    \fill[bsWater!16!white]
      (-1.9, 2.05) .. controls (0.5, 3.85) and (3.5, 3.85) .. (5.9, 2.05)
      -- (5.9, 0.75) .. controls (3.5, 2.40) and (0.5, 2.40) .. (-1.9, 0.75) -- cycle;

    % Banco de sedimentos na margem convexa (interna)
    \fill[bsGold!28!white]
      (-0.6, 1.25) .. controls (1.2, 2.30) and (2.8, 2.30) .. (4.6, 1.25)
      .. controls (3.2, 1.90) and (0.8, 1.90) .. (-0.6, 1.25) -- cycle;

    % Contorno das margens
    \draw[bsEarth, line width=1.6pt]
      (-1.9, 2.05) .. controls (0.5, 3.85) and (3.5, 3.85) .. (5.9, 2.05);
    \draw[bsGoldDark, line width=1.4pt]
      (-1.9, 0.75) .. controls (0.5, 2.40) and (3.5, 2.40) .. (5.9, 0.75);

    % Talvegue (deslocado para a margem côncava)
    \draw[bsNavy, dashed, line width=1.1pt]
      (-1.8, 1.45) .. controls (0.6, 3.38) and (3.4, 3.38) .. (5.8, 1.45);

    % Vetores de velocidade
    \draw[->, >=Stealth, bsWater, line width=1.35pt] (-1.1, 2.15) -- (0.1, 2.85);
    \draw[->, >=Stealth, bsWater, line width=1.6pt]  (0.7, 3.10) -- (2.5, 3.28);
    \draw[->, >=Stealth, bsWater, line width=1.5pt]  (2.8, 3.18) -- (4.3, 2.65);
    \draw[->, >=Stealth, bsWater!75!white, line width=0.95pt] (1.3, 2.45) -- (2.1, 2.52);

    % Indicadores de solapamento na margem côncava
    \foreach \x/\y in {1.1/3.55, 2.0/3.66, 2.9/3.55} {
      \draw[->, >=Stealth, bsAlert, line width=1.1pt] (\x, \y-0.24) -- (\x, \y+0.12);
    }

    \node[font=\scriptsize\bfseries, text=bsGoldDark, align=center, fill=bgGold, draw=bsGold, rounded corners=3pt, inner xsep=5pt, inner ysep=3pt]
      at (2.05, 0.92) {Margem Convexa (Interna): Baixa Energia \& Deposição de Sedimentos};
    \node[font=\scriptsize\itshape, text=bsNavy] at (4.75, 3.05) {Talvegue};
  \end{scope}

  % Linha divisória interna
  \draw[bsMuted!40!white, dashed, line width=0.8pt] (-2.35, 0.85) -- (6.55, 0.85);

  % Seção Transversal Assimétrica (A-A')
  \begin{scope}[shift={(0.1, -2.25)}]
    \node[font=\scriptsize\bfseries, text=bsMuted, anchor=west] at (-2.45, 2.70) {SEÇÃO TRANSVERSAL};

    % Água na seção
    \fill[bsWater!20!white]
      (-0.6, 1.75) -- (4.7, 1.75) -- (4.7, 1.55)
      .. controls (3.1, 1.45) and (1.8, 0.75) .. (0.6, 0.45)
      .. controls (0.0, 0.35) and (-0.35, 1.1) .. (-0.6, 1.75) -- cycle;

    % Nível d'água
    \draw[bsWater, dashed, line width=1pt] (-0.8, 1.75) -- (5.0, 1.75);
    \node[font=\scriptsize\bfseries, text=bsWater, fill=white, rounded corners=2pt, inner sep=1.5pt] at (3.75, 1.98) {N.A.};

    % Perfil do leito e margens
    \fill[bsEarth!14!white]
      (-2.1, 2.20) -- (-0.9, 2.20) -- (-0.55, 1.65)
      .. controls (-0.3, 1.0) and (0.0, 0.35) .. (0.6, 0.45)
      .. controls (1.8, 0.75) and (3.2, 1.50) .. (4.9, 1.65)
      -- (5.8, 1.95) -- (5.8, 0.0) -- (-2.1, 0.0) -- cycle;

    \draw[bsEarth, line width=1.6pt]
      (-2.1, 2.20) -- (-0.9, 2.20) -- (-0.55, 1.65)
      .. controls (-0.3, 1.0) and (0.0, 0.35) .. (0.6, 0.45)
      .. controls (1.8, 0.75) and (3.2, 1.50) .. (4.9, 1.65) -- (5.8, 1.95);

    % Depósito na barra em pontal
    \fill[bsGold!35!white]
      (2.7, 1.32) .. controls (3.5, 1.55) and (4.2, 1.62) .. (4.8, 1.64)
      .. controls (4.0, 1.35) and (3.3, 1.22) .. (2.7, 1.32) -- cycle;

    % Tensão cisalhante e poço de erosão
    \draw[->, >=Stealth, bsAlert, line width=1.3pt] (0.55, 0.88) -- (-0.15, 0.68);
    \node[font=\scriptsize\bfseries, text=bsAlert, align=center] at (-1.15, 0.65)
      {Fosso de\\solapamento ($\tau_0$)};
    \node[font=\scriptsize\bfseries, text=bsGoldDark, align=center] at (4.15, 1.02)
      {Barra de\\sedimentos};
    \node[font=\scriptsize\bfseries, text=bsNavy] at (1.55, 1.28) {Núcleo de alta $V$};
  \end{scope}

  % ============================================================
  % PAINEL DIREITO: ZONAS ECOLÓGICAS E HIDRÁULICAS NA MARGEM
  % ============================================================
  \node[bsTint=bsForest, minimum width=6.15cm, minimum height=8.85cm, anchor=north west] (rightpanel) at (7.3, 6.10) {};
  \node[bsBadge=bsForest] at (rightpanel.north) {ZONAS DE ATUAÇÃO NA MARGEM};

  % Barra lateral de cota hidráulica
  \draw[bsWater, line width=1.6pt]  (7.80, -2.35) -- (7.80, 0.35);
  \draw[bsTeal, line width=1.6pt]   (7.80, 0.35)  -- (7.80, 3.00);
  \draw[bsForest, line width=1.6pt] (7.80, 3.00)  -- (7.80, 5.45);
  \node[font=\tiny\bfseries, text=bsMuted, rotate=90] at (7.55, 1.55) {COTA HIDRÁULICA CRESCENTE $\rightarrow$};

  % Zona 3: Terrestre / APP (Topo)
  \node[bsBox=bsForest, text width=4.65cm, minimum height=2.25cm, anchor=north west] (z3) at (8.15, 5.45) {
    \vspace{0.18cm}
    \textbf{\color{bsForest}Zona Terrestre / APP}\\[2pt]
    \scriptsize
    \textbf{Mata Ciliar Nativa}\\
    Espécies arbóreas e arbustivas de raízes profundas; interceptação pluvial, ancoragem do topo e conectividade ecológica.
  };
  \node[bsBadge=bsForest, font=\tiny\bfseries] at (z3.north) {ACIMA DA CHEIA MÁXIMA};

  % Zona 2: Transição / Nível Médio
  \node[bsBox=bsTeal, text width=4.65cm, minimum height=2.25cm, anchor=north west] (z2) at (8.15, 2.75) {
    \vspace{0.18cm}
    \textbf{\color{bsTeal}Zona de Transição / Nível Médio}\\[2pt]
    \scriptsize
    \textbf{Feixes \& Gabião Vivo}\\
    Estacas vivas, esteiras de ramagem e gabiões vegetados; dissipação em cheias sazonais e retenção de finos.
  };
  \node[bsBadge=bsTeal, font=\tiny\bfseries] at (z2.north) {VARIAÇÃO SAZONAL DO N.A.};

  % Zona 1: Submersa / Pé (Base)
  \node[bsBox=bsEarth, text width=4.65cm, minimum height=2.25cm, anchor=north west] (z1) at (8.15, 0.05) {
    \vspace{0.18cm}
    \textbf{\color{bsEarth}Zona Submersa / Pé da Margem}\\[2pt]
    \scriptsize
    \textbf{Enrocamento \& Espigão}\\
    Proteção estrutural rígida ou semirrígida contra tensão cisalhante ($\tau_0$) contínua e solapamento basal.
  };
  \node[bsBadge=bsEarth, font=\tiny\bfseries] at (z1.north) {SUBMERSA PERMANENTE};

\end{tikzpicture}
\end{document}
